% !PRINTED 110
\weberpage{110}
%<<CONT
und ebenso erhalten wir

\weq{14}{(\varrho^2, x_0) = (-\varepsilon^2)^\lambda \varrho_2 (-\sqrt{-3})^u q_1^{\sigma_1} q_2^{\sigma_2} \ldots p_1^{-\tau_1} p_2^{-\tau_2} \ldots H_2.} % [sic: Weber prints exponent u; the conjugate of (13) requires \mu.]

Weil das Product dieser beiden Ausdrücke rational sein muss, so ergiebt sich
noch zwischen den Einheitswurzeln $\varrho, \varrho_1, \varrho_2$ die Relation

\[
\varrho^\lambda \varrho_1 \varrho_2 = 1.
\]

Setzen wir

\weq{15}{x_0 + x_1 + x_2 = A,}

so ist auch $A$ eine rationale Zahl, und aus (13), (14), (15) folgt, dass

\weq{16}{\begin{aligned}
x_0 &= 1/3 \left[A + (\varrho, x_0) + (\varrho^2, x_0)\right] \\
x_1 &= 1/3 \left[A + \varrho^2 (\varrho, x_0) + \varrho (\varrho^2, x_0)\right] \\
x_2 &= 1/3 \left[A + \varrho (\varrho, x_0) + \varrho^2 (\varrho^2, x_0)\right]
\end{aligned}}

Kreistheilungzahlen sind. % [sic] printed \textquotedbl{}Kreistheilungzahlen\textquotedbl{}, elsewhere on the page \textquotedbl{}Kreistheilungszahlen\textquotedbl{}
Aus den Formeln (13), (14) können wir die Zusammensetzung dieser Zahlen aus den
Perioden $\eta_1, \eta_2, \ldots$ ersehen. Ein näheres Eingehen auf diesen
Gegenstand ist aber nicht mehr erforderlich, weil wir schon im \S.\,23 alle
Kreistheilungszahlen, die cubischen Gleichungen genügen, vollständig gebildet
haben.

\noindent\hrulefill

\begin{center}
\S.\,26.

\gesperrt{Biquadratische Abel'sche Gleichungen.}
\end{center}

Ganz ähnlich kann der Beweis des entsprechenden Satzes für die biquadratischen
\gesperrt{Abel}'schen Gleichungen geführt werden. Nur sind hier zwei Fälle zu
unterscheiden, nämlich Gleichungen mit nicht cyklischer Gruppe, und
Gleichungen mit cyklischer Gruppe.

Für die nicht cyklischen \gesperrt{Abel}'schen Gleichungen mit den Wurzeln $x_0, x_1, x_2,
x_3$ ist die Gruppe:

\weq{1}{1, \ \ (0,1)(2,3), \ \ (0,2)(1,3), \ \ (0,3)(1,2),}

und es sind also die drei Quadrate

\weq{2}{\begin{aligned}
(x_0+x_1-x_2-x_3)^2 &= a \\
(x_0-x_1+x_2-x_3)^2 &= b \\
(x_0-x_1-x_2+x_3)^2 &= c
\end{aligned}}

und das Product

\weq{3}{(x_0+x_1-x_2-x_3)(x_0-x_1+x_2-x_3)(x_0-x_1-x_2+x_3) = c,} % [sic: Weber prints RHS c, reusing c from the third square in (2); this rational product requires a distinct symbol.]
%CONT>>
